% Generated draft: requires mathematical and source review.
\documentclass[11pt]{article}
\usepackage[margin=25mm]{geometry}
\usepackage{fontspec}
\setmainfont{Latin Modern Roman}
\usepackage{amsmath,amssymb,mathtools,unicode-math}
\setmathfont{Latin Modern Math}
\usepackage{enumitem,needspace,xurl,hyperref}
\setlength{\parindent}{0pt}
\setlength{\parskip}{5pt}
\setlength{\emergencystretch}{3em}
\newcommand{\N}{\mathbb{N}}
\newcommand{\Z}{\mathbb{Z}}
\newcommand{\Q}{\mathbb{Q}}
\newcommand{\R}{\mathbb{R}}
\newcommand{\C}{\mathbb{C}}
\newcommand{\F}{\mathbb{F}}
\newcommand{\A}{\mathbb{A}}
\newcommand{\PP}{\mathbb{P}}
\newcommand{\E}{\mathbb{E}}
\newcommand{\eps}{\varepsilon}
\newcommand{\dd}{\,\mathrm{d}}
\begin{document}
\section*{UnsolvedMath ID: 30002317}
\textbf{Source problem code:} catalogue entry
\subsection{Definitions and mathematical statement}
The supplied record leaves several essential conventions unspecified. The following gives a conditional formulation and identifies the unresolved parts; it does not supply missing hypotheses.

For a ground field \(k\) and an integer \(n\geq 1\), write \(\mathbb A_k^n\) for affine \(n\)-space and use the standard notation
\[
\operatorname{Aff}_n(k)=k^n\rtimes\operatorname{GL}_n(k),\qquad
\operatorname{Cr}_n(k)=\operatorname{Bir}_k(\mathbb P_k^n).
\]
Here an element of \(\operatorname{Aff}_n(k)\) acts by \(x\mapsto Ax+b\), and \(\operatorname{Aut}_k(\mathbb A_k^n)\) denotes the group of regular \(k\)-automorphisms of affine space. Affine automorphisms are viewed as elements of the Cremona group through the standard affine chart. The source does not specify \(k\), its characteristic, whether it is algebraically closed, or the admissible dimension range. Thus this notation does not select the intended ground-field hypotheses or settle whether dimension zero is included.

The notation \(\operatorname{Aut}^*(\mathbb A^n)\) and the meaning of a \emph{torsion prime} of these groups are not defined in the supplied record. Retain \(\operatorname{Tors}(G)\) as the source's unspecified set of torsion primes. In particular, no identification with primes occurring as orders of group elements is made.

Subject to these unresolved conventions, the questions are as follows.
\begin{enumerate}
\item For which admissible dimensions \(n\) do there exist finite subgroups \(H_1,H_2\leq\operatorname{Aff}_n(k)\) such that
\[
H_1\cong H_2
\quad\text{and}\quad
\nexists\,g\in\operatorname{Cr}_n(k)\text{ with }gH_1g^{-1}=H_2?
\]
The supplied wording does not settle whether the intended question concerns a prescribed dimension or existence in some dimension.
\item For each admissible \(n\), are all three sets
\[
\operatorname{Tors}(\operatorname{Cr}_n(k)),\qquad
\operatorname{Tors}(\operatorname{Aut}_k(\mathbb A_k^n)),\qquad
\operatorname{Tors}(\operatorname{Aut}^*(\mathbb A_k^n))
\]
finite? The third expression retains the undefined starred notation from the source.
\item Determine whether there is an admissible dimension \(n\) for which
\[
7\in
\operatorname{Tors}(\operatorname{Cr}_n(k))
\cup\operatorname{Tors}(\operatorname{Aut}_k(\mathbb A_k^n))
\cup\operatorname{Tors}(\operatorname{Aut}^*(\mathbb A_k^n)),
\]
and, if so, determine the least such \(n\). This follows the supplied wording ``one of them''; the record does not explicitly request separate minima for the three families.
\item Determine which analogous or modified Jordan-decomposition properties for connected reductive algebraic groups hold for these three groups. The record does not enumerate the properties to be tested or define the intended notions of semisimple and unipotent elements in these groups. Consequently, this part cannot yet be expressed as a definite collection of mathematical assertions.
\end{enumerate}
\subsection{Short English statement}
The problem asks whether isomorphic finite affine transformation groups can fail to be conjugate by a Cremona transformation, whether three specified torsion-prime sets are finite, when the prime \(7\) first occurs in their union, and which Jordan-decomposition properties extend to these groups. The ground field, torsion-prime convention, starred automorphism group, and precise Jordan-decomposition targets require clarification. The supplied status is \emph{partially solved}; this is a curation label, not an independently verified conclusion.
\subsection{Sources}
\begin{itemize}
\item Supplied local UnsolvedMath registry record, permanent ID \(30002317\), title \emph{Finite Subgroups and Torsion Primes of Cremona Groups}, source problem code \texttt{catalogue entry}, in the Oberwolfach Reports Open Problems collection. No individual proposer is supplied; the proposed year is recorded as 2013.
\item Supplied report citation: \emph{Algebraic Groups (2014)}, Oberwolfach Reports, DOI \texttt{10.4171/owr/2013/17}, pp.\ 1025--1085. \url{https://doi.org/10.4171/owr/2013/17}. The supplied title year is 2014, whereas the DOI and proposed-year field contain 2013; this discrepancy is unresolved.
\item Supplied literature reference: Popov, \emph{Tori in the Cremona groups}. \url{https://www.math.uni-bielefeld.de/LAG/man/474.pdf}. No further bibliographic details or precise supporting passages are supplied.
\item Supplied literature reference: \emph{Oberwolfach algebraic-groups report}. \url{https://doi.org/10.4171/owr/2013/17}.
\item Literature assessment dated exactly \texttt{2026-08-21}: the supplied curation labels the problem partially solved and reports known nonconjugate finite-subgroup phenomena and some Cremona torsion information, while describing the affine-conjugacy, automorphism-group torsion, and minimal-prime questions as open. This is supplied dated curation derived from the OWR material, not independently verified here. The background explicitly requests independent verification of time-sensitive status and cited claims.
\item The registry labels the statement \texttt{corrected\_verified} and reports that it was checked against the official Oberwolfach report, with source hypotheses and notation restored. That verification claim is supplied metadata, not verification performed for this draft; the supplied text still omits essential conventions.
\end{itemize}
\end{document}
